\documentclass[11pt]{article}

\usepackage[margin=1in]{geometry}
\usepackage{amsmath, amssymb}
\usepackage[T1]{fontenc}
\usepackage[utf8]{inputenc}
\usepackage{lmodern}
\usepackage{enumitem}
\usepackage{hyperref}

\setlist[itemize]{topsep=4pt, itemsep=2pt, parsep=0pt}

\title{6.8610 Spring 2026 Lecture 1 Notes: Introduction to NLP and Language Models}
\author{}
\date{February 3, 2026}

\begin{document}
\maketitle

\section{Introduction to NLP and Language Models}

\subsection{Motivation}
\textbf{Definition: Natural Language Processing (NLP)}

NLP studies how to make computers process, approximate understanding of, and leverage human language data. In this course, the emphasis is on \textbf{written text}, although many of the same modeling ideas extend to speech and signed languages.

\begin{itemize}
    \item Modern society produces enormous quantities of text, from web pages and news articles to search queries and social media posts.
    \item If computers can model human language well, they can support search, translation, summarization, question answering, assistants, and many other tools that directly help people.
    \item The lecture emphasizes that recent progress has made NLP visible to the general public, but current systems are still imperfect and should not be treated as infallible or human.
    \item A central modern insight is that many NLP tasks can be reframed as \textbf{language modeling}: learning probability distributions over strings, or conditional distributions that generate one text from another.
\end{itemize}

\subsection{Key Concepts}
\textbf{Definition: Language Model}

A language model assigns probabilities to strings:
\[
p_\theta(x_{1:T})
\]
where \(x_{1:T} = (x_1, x_2, \dots, x_T)\) is a token sequence and \(\theta\) denotes model parameters.

\begin{itemize}
    \item \textbf{Human language is special and difficult to model.}
    \begin{itemize}
        \item It is discrete: language is built from compositional units such as tokens, morphemes, or words.
        \item It is acquired remarkably naturally by humans.
        \item It is almost infinitely expressive: the same meaning can be phrased many ways, and new meanings can be composed productively.
        \item It is socially grounded, nuanced, and constantly evolving.
    \end{itemize}
    \item \textbf{Surface form and meaning are not the same.} The same underlying concept can be expressed across different encodings and languages.
    \item \textbf{Language is ambiguous.} Meaning depends on context, world knowledge, discourse, syntax, and even stylistic choices such as punctuation or capitalization.
    \item \textbf{Representation and modeling are the two cornerstones of NLP.}
    \begin{itemize}
        \item \textbf{Representation}: how symbolic language is converted into a form a computer can process.
        \item \textbf{Modeling}: how those representations are used to solve a task or estimate a probability distribution.
    \end{itemize}
    \item \textbf{NLP historically consisted of many specialized tasks.} These include syntax-focused tasks (morphology, part-of-speech tagging, parsing), discourse tasks (summarization, coreference), and semantic tasks (sentiment, named entity recognition (NER), question answering (QA), machine translation, language modeling).
    \item \textbf{Modern NLP increasingly unifies tasks under language modeling.} If a model can produce or condition on text well enough to follow instructions, many traditional NLP tasks become special cases of text generation or conditional prediction.
\end{itemize}

\subsection{Core Models and Equations}
\textbf{Task framing.} Many NLP problems can be written as learning a function from an input \(X\) to an output \(Y\):
\[
Y \approx f_\theta(X)
\]
For example, in machine translation, \(X\) may be an English sentence and \(Y\) its Spanish translation.

\textbf{Maximum-likelihood language modeling.} Given a training corpus \(x_{1:T}\), the goal is to learn parameters that maximize the probability of the observed data:
\[
\hat{\theta} = \arg\max_{\theta} p_\theta(x_{1:T})
\]
Often, the training corpus is viewed as one very long token sequence obtained by concatenating many documents.

\textbf{Autoregressive factorization.} By the chain rule,
\[
p_\theta(x_{1:T}) = \prod_{t=1}^{T} p_\theta(x_t \mid x_{1:t-1})
\]
This left-to-right factorization is the dominant paradigm in modern large language models (LLMs).

\textbf{Alternative factorizations.} The lecture notes that the joint distribution can also be factorized in other ways, for example:
\[
p_\theta(x_{1:T}) = \prod_{t=1}^{T} p_\theta(x_t \mid x_{t+1:T})
\]
for right-to-left generation, or via latent-variable models such as hidden Markov model (HMM)-style models or diffusion-based formulations.

\textbf{Conditional language modeling.} Many downstream tasks can be reframed as learning:
\[
p_\theta(y_{1:m} \mid x_{1:n}) = \prod_{t=1}^{m} p_\theta(y_t \mid y_{1:t-1}, x_{1:n})
\]
This viewpoint underlies sequence-to-sequence learning for tasks such as translation and summarization.

\textbf{Naive scalar parameterization.} If we explicitly stored the next-token distribution for every possible context, then with vocabulary size \(V\) and maximum length \(T\), the number of scalar probabilities grows as
\[
\sum_{t=1}^{T} V^t
\]
which is exponential in sequence length. This is why direct table lookup does not scale and function approximation is required.

\textbf{Perplexity (PPL).} A standard intrinsic evaluation metric for language models is
\[
\operatorname{PPL}(x_{1:T}) = \exp\left(-\frac{1}{T}\log p_\theta(x_{1:T})\right)
\]
Using the autoregressive factorization,
\[
\operatorname{PPL}(x_{1:T}) = \exp\left(-\frac{1}{T}\sum_{t=1}^{T}\log p_\theta(x_t \mid x_{1:t-1})\right)
\]
Equivalent forms include
\[
\operatorname{PPL}(x_{1:T}) = p_\theta(x_{1:T})^{-1/T}
\]
Perplexity can be interpreted as the model's \textbf{average branching factor}; for a uniform distribution over a vocabulary of size \(V\), perplexity equals \(V\).

\subsection{Algorithms / Architectures}
\textbf{Historical modeling progression}
\begin{itemize}
    \item \textbf{1960s}: pattern matching and rule-based systems.
    \item \textbf{1970s--1980s}: linguistically rich, logic-driven systems.
    \item \textbf{1990s--2000s}: statistical machine learning becomes central.
    \item \textbf{2010s--2020s}: deep neural networks dominate NLP.
    \item \textbf{2020s onward}: large language models provide a single flexible interface for many tasks.
\end{itemize}

\textbf{Generic supervised NLP pipeline}
\begin{itemize}
    \item Start with input data \(X\) and desired outputs \(Y\).
    \item Choose a model family \(f_\theta\): rule-based, conditional random field (CRF), HMM, graphical model, neural network, or another structured predictor.
    \item Fit parameters to data so that the model maps \(X\) to useful outputs \(Y\).
    \item Evaluate according to the application: accuracy, robustness, fairness, latency, or human usefulness.
\end{itemize}

\textbf{Autoregressive generation procedure}
\begin{itemize}
    \item Begin with a start-of-sequence context.
    \item Predict a distribution over the next token.
    \item Select or sample the next token.
    \item Append it to the context.
    \item Repeat until an end-of-sequence token is produced.
\end{itemize}

\textbf{Why function approximation matters}
\begin{itemize}
    \item A literal count table over all contexts is exponentially large.
    \item Neural networks replace explicit memorization with shared parameters.
    \item Shared parameters allow generalization across related contexts, tokens, and tasks.
    \item In the deep-learning era, many systems learn features automatically instead of relying on manual feature engineering.
\end{itemize}

\subsection{Important Intuitions from the Professor}
\begin{itemize}
    \item \textbf{Do not anthropomorphize language models.} Fluency is not the same as sentience or grounded understanding.
    \item \textbf{High-stakes decisions still require human oversight.} The lecture explicitly warns against treating model outputs as authoritative for serious matters such as medicine or major personal decisions.
    \item \textbf{``Understanding'' may be the wrong operational target.} In practice, NLP systems often succeed by \emph{approximating understanding} well enough for a specific objective.
    \item \textbf{Labeled data is scarce and valuable.} Unlabeled text is abundant, but high-quality task annotations are expensive and remain a bottleneck.
    \item \textbf{Language evolves.} New words and usages such as ``COVID-19'' or ``Brexit'' force models to handle novelty and distribution shift.
    \item \textbf{The right scientific habit is to inspect data and outputs carefully.} The professor emphasized that strong empirical work requires checking simple examples and failure cases rather than only reporting aggregate scores.
    \item \textbf{Representation choices matter as much as model choices.} Tokenization and input formatting are not bookkeeping details; they shape what the model can learn.
    \item \textbf{Modern NLP has consolidated around language modeling.} Earlier systems were task-specific; now a sufficiently strong language model that can follow instructions can support many text-processing tasks through a shared interface.
\end{itemize}

\subsection{Examples}
\begin{itemize}
    \item \textbf{Ambiguity examples from the slides}
    \begin{itemize}
        \item ``Red tape holds up new bridges.'' This can describe either literal support or bureaucratic delay.
        \item ``Tesla crashed today.'' The subject could be a car, a stock, or a person named Tesla.
        \item ``She made him duck.'' The sentence could mean forcing someone to lower their head or cooking duck for him.
        \item ``Yes'' vs.\ ``Yes.'' vs.\ ``YES!'' show that punctuation and capitalization alter pragmatic meaning.
    \end{itemize}
    \item \textbf{Voice-assistant intent extraction}
    \begin{itemize}
        \item In ``Alexa, play Drivers License by Olivia Rodrigo,'' a system must infer the \textbf{intent} (play music) and extract slots such as song title and artist.
        \item The professor gave an edge case: ``Play David Bowie by Phish,'' where both phrases can refer to artists, but the intended parse is song-by-artist.
    \end{itemize}
    \item \textbf{Machine translation}
    \begin{itemize}
        \item ``El perro marr\'on'' \(\rightarrow\) ``The brown dog'' illustrates that equivalent meanings can require different word orders across languages.
    \end{itemize}
    \item \textbf{Autoregressive probability}
    \begin{itemize}
        \item For ``I am a cat,'' the model computes the joint probability as a product of next-token probabilities conditioned on the prefix.
    \end{itemize}
    \item \textbf{Rule-based sentiment failure}
    \begin{itemize}
        \item A rule such as ``if the review contains \emph{incredible}, predict five stars'' fails on inputs like ``incredibly bad'' or more compositional negations and discourse twists.
    \end{itemize}
\end{itemize}

\subsection{Common Pitfalls / Limitations}
\begin{itemize}
    \item \textbf{Rule brittleness}: handcrafted rules break under negation, irony, compositionality, and domain shift.
    \item \textbf{Surface-form confusion}: identical strings can have different meanings, and different strings can express the same meaning.
    \item \textbf{Data mismatch}: systems that work on one corpus may fail badly on another because vocabulary, style, or entity distributions change.
    \item \textbf{Exponential context growth}: naive scalar next-token tables do not scale to realistic vocabularies and contexts.
    \item \textbf{Overtrust}: fluent outputs can conceal factual errors, bias, or harmful advice.
    \item \textbf{Metric myopia}: a ``good'' model depends on the use case; accuracy alone may ignore fairness, robustness, latency, or usability.
\end{itemize}

\subsection{Connections to Other Lectures}
As the opening lecture, this class session establishes the conceptual frame for the rest of 6.8610 rather than building on earlier course lectures.

\begin{itemize}
    \item The distinction between \textbf{representation} and \textbf{modeling} foreshadows later lectures on tokenization, architectures, pretraining, post-training, and ranking.
    \item The shift from many specialized models to unified \textbf{language modeling} motivates why later lectures focus on log-linear models, recurrent neural networks (RNNs), sequence-to-sequence models, attention, Transformers, and modern LLM training.
    \item The introduction of \textbf{perplexity} prepares for later discussions of evaluation, scaling, and why intrinsic metrics do not fully capture downstream usefulness.
    \item The lecture also connects directly to standard machine-learning prerequisites: maximum likelihood, supervised learning, generalization, and parameterized function approximation.
\end{itemize}

\subsection{Exam-Relevant Summary}
\begin{itemize}
    \item \textbf{NLP definition}: making computers process and leverage human language, with this course focusing on written text.
    \item \textbf{Why NLP is hard}: ambiguity, compositionality, pragmatics, social context, and language change.
    \item \textbf{Two cornerstones}: representation and modeling.
    \item \textbf{Classical NLP task families}: syntax, discourse, and semantics.
    \item \textbf{Modern unification}: many tasks can be reframed as unconditional or conditional language modeling.
    \item \textbf{Language model}: a probability distribution over token sequences.
    \item \textbf{Autoregressive factorization}:
    \[
    p_\theta(x_{1:T}) = \prod_{t=1}^{T} p_\theta(x_t \mid x_{1:t-1})
    \]
    \item \textbf{Training objective}: maximum likelihood on a large text corpus.
    \item \textbf{Naive scalar parameterization is exponential} in context length, motivating learned function approximation.
    \item \textbf{Perplexity} is the exponential of average negative log-likelihood and can be interpreted as average branching factor.
    \item \textbf{Historical progression}: rules \(\rightarrow\) linguistic/statistical models \(\rightarrow\) deep learning \(\rightarrow\) LLMs.
    \item \textbf{Professor emphasis}: do not confuse fluent text with true understanding; inspect data and model failures carefully.
\end{itemize}

\end{document}
